% This is samplepaper.tex, a sample chapter demonstrating the
% LLNCS macro package for Springer Computer Science proceedings;
% Version 2.21 of 2022/01/12
%
\documentclass[runningheads]{llncs}
%
\usepackage[T1]{fontenc}
% T1 fonts will be used to generate the final print and online PDFs,
% so please use T1 fonts in your manuscript whenever possible.
% Other font encondings may result in incorrect characters.
%
\usepackage{graphicx}
% Used for displaying a sample figure. If possible, figure files should
% be included in EPS format.
%
%
% Added by authors
\usepackage{csquotes}
\usepackage{amssymb}
\usepackage{amsmath}
\usepackage{todonotes}
\usepackage{hyperref}
\usetikzlibrary{arrows.meta,calc}

% If you use the hyperref package, please uncomment the following two lines
% to display URLs in blue roman font according to Springer's eBook style:
\usepackage{color}
\renewcommand\UrlFont{\color{blue}\rmfamily}
\urlstyle{rm}
%
\begin{document}
%
\title{Relaxing Global Classifier Robustness on Semimetric Spaces}
%
%\titlerunning{Global Robustness: A Feasibility Threshold and its Relaxation}
% If the paper title is too long for the running head, you can set
% an abbreviated paper title here
%
\author{Andrei Aleksandrov\inst{1,2}\orcidID{0000-0002-4717-4206} \and
Kim Völlinger\inst{2}\orcidID{0000-0002-8988-0053}}
%
\authorrunning{A. Aleksandrov and K. Völlinger}
% First names are abbreviated in the running head.
% If there are more than two authors, 'et al.' is used.
%
\institute{Fraunhofer Institute for Open Communication Systems, Berlin, Germany \\
\email{andrei.aleksandrov@fokus.fraunhofer.de} \and
Technische Universität Berlin, Berlin, Germany \\
\email{voellinger@tu-berlin.de}}
%
\maketitle              % typeset the header of the contribution
%
\begin{abstract}
A classifier is $\epsilon$-robust if it assigns the same label to any two inputs within distance $\epsilon$.
On continuous domains, every $\epsilon$-robust classifier is constant, motivating existing local and probabilistic relaxations.
We introduce $(\mu,\epsilon)$-robustness, which requires every input to lie within distance $\mu$ of a same-label point at which the classifier is locally $\epsilon$-robust. This notion recovers $\epsilon$-robustness when $\mu=0$ and, unlike probabilistic relaxations, imposes explicit coverage of the domain by robust points. We study robustness over semimetric spaces, with a focus on the finite, totally ordered domains serving as a model of finite-precision arithmetic. We analyze the feasibility of $(\mu,\epsilon)$-robustness independently of accuracy. In particular, we prove a covering characterization of robustness, derive necessary and sufficient existence criteria for robust non-constant classifiers, and establish upper bounds on $\epsilon$ and $\mu$. 

\keywords{Classifier \and Robustness \and Computer Arithmetic.}
\end{abstract}

\section{Introduction}

Machine learning classifiers can change their prediction under small input perturbations, raising concerns about safety, security, and ethics. 
A central notion for studying this phenomenon is \emph{global robustness}, which requires that similar inputs receive identical (or similar) outputs~\cite{casadio2022neural}. 
Similarity is captured by a distance function, making the framework remarkably general: in computer vision it led to weaker notions, since similarity based on pixel intensities propagates by transitivity; but
in fairness, where the distance encodes context-dependent similarity of individuals under ethical constraints (e.g., equally qualified for a job)~\cite{fairness-landscape,kusner2017counterfactual,ranking-fairness}, and in natural language processing, syntactic and semantic similarity is encoded under meaningful linguistic variation~\cite{levy2025statistical} remains central.
Importantly, the notion of similarity may \emph{not} be metric, violating the triangle inequality~\cite{garcia2019learning,similarity-guide}.
Motivated by this observation, we study robustness over arbitrary \emph{semimetric} spaces.
A common observation is that exact global $\epsilon$-robustness is essentially unsatisfiable on connected domains such as $\mathbb{R}^n$: overlapping $\epsilon$-balls form chains that connect any two inputs, and force a robust classifier to be constant. 
Existing work has addressed this limitation by weakening the guarantee: 
local robustness concerns only individual inputs~\cite{li2024certifying}; confidence-based notions restrict robustness to high-confidence
inputs~\cite{athavale2024verifying,leino2021globally}; probabilistic global robustness~\cite{mangal2019robustness,levy2023groma}.
However, probabilistic notions constrain only the density of counterexamples, not their geometric structure (see Fig.~\ref{fig:grid}).

To address this limitation, we introduce $(\mu,\epsilon)$-robustness:
Every input must lie within distance $\mu$ of a same-label point at which the classifier is locally $\epsilon$-robust. 
The definition recovers $\epsilon$-robustness at $\mu = 0$ and, unlike probabilistic
relaxations, imposes an explicit coverage \emph{structure} on non-robust points.
Two of our main results, a covering characterization of robustness and a necessary-and-sufficient criterion for the
existence of a non-constant robust classifier, hold over arbitrary semimetric spaces,
providing a domain-independent foundation.
Further, the impossibility result above is a consequence of \emph{connected
continuous} domains, whereas deployed classifiers compute over finite-precision
arithmetic. In fact, verification methods reasoning over reals can be unsound 
due to their abstraction~\cite{jia2021exploiting,cordeiro2025neural}. 
Since deployed systems operate over finite representable values rather than continuous domains, robustness becomes a quantitative question.
We show that the existence of a non-constant robust classifier is determined exactly 
by the relationship between $\epsilon$ ($\mu$) and the gaps between consecutive representable values. 

\paragraph{Contributions.}
Studying $(\mu,\epsilon)$-robustness gives us domain-independent structural results and quantitative guarantees for finite-precision arithmetic:
\begin{enumerate}
\item \textbf{Covering characterization.}
We introduce a covering characterization of $(\mu,\epsilon)$-robustness (Section~\ref{sec:partitions-covering}), extending the classical connection between $\epsilon$-robustness and separability. In the distribution-free setting, we establish separability as a necessary condition for global robustness~\cite{yang2020closer,mangal2023certifying}.

\item \textbf{Existence criteria.}
We derive necessary and sufficient conditions for the existence of a non-constant robust classifier (Section~\ref{sec:existence}), both for arbitrary semimetric spaces and for finite totally ordered domains.

\item \textbf{Precision-dependent bounds.}
We obtain bounds on $\epsilon$ and $\mu$ in terms of precision (ULP) metrics (Section~\ref{sec:upper-bounds}). These bounds identify the largest representable gap as the key quantity governing robustness, explain empirically observed robustness gains under coarser quantization~\cite{xiao2023robustmq,shrestha2024properties,duncan2020relative}, and extend to some partially ordered domains.
\end{enumerate}
Our analysis characterizes \emph{when} robust classifiers exist. 
Whether an \emph{accurate} robust classifier exists, and thus the robustness-accuracy
tradeoff~\cite{zhang2019theoretically}, is orthogonal to our focus. 
Along the way, our framework provides a unified, non-probabilistic perspective on several results 
concerning global robustness that have previously appeared in disparate parts of the literature.

%\begin{enumerate}
%    \item In Section~\ref{sec:new-robustness}, we define $(\mu,\epsilon)$-robustness and discuss alternative formulations of global classifier robustness.
%    \item In Section~\ref{sec:partitions-covering}, we formalize $(\mu,\epsilon)$-robustness on induced by classifiers partitions and show its equivalence to a covering property.
%    \item In Section~\ref{sec:existence}, we establish necessary and sufficient existence criteria for a non-constant robust classifier in general case and on finite totally ordered semimetric spaces, generalizing a line of finite precision numbers.
%    \item In Section~\ref{sec:upper-bounds}, we present upper bounds on $\epsilon$ and $\mu$ for non-constant classifiers on finite totally ordered semimetric spaces as well as their products.
%\end{enumerate}
\section{$(\mu,\epsilon)$-Robustness}
\label{sec:new-robustness}

We define classifier robustness in its usual form, survey related work containing other approaches to defining global robustness and highlight differences to our proposed formulation.

\paragraph{Local Robustness.} Local classifier robustness is an intuitive widely accepted in the literature property, which we formulate over semimetric spaces. We use the notation $[n] := \{1, 2, \dots,n\} \subset \mathbb{N}$.
% lax begin Lax544225.Classifier
\begin{definition}[Classifier]
\label{def:classifier}
Let $S$ be a set of \enquote{inputs}  and  $[n]$ a set of \enquote{classes}. 
A \emph{classifier} is a total, surjective function:
$c: S \rightarrow [n]$. 
\end{definition}
% lax end

Surjectivity ensures that each class is represented by at least one input, reflecting the common requirement in machine learning that all labels occur in the domain and helping to avoid systematic bias toward missing classes.
% lax begin Lax544225.Semimetric
\begin{definition}[Semimetric Space]
\label{def:distance}
    Let $(S, d)$ be a pair of a nonempty set and a function $d: S \times S \rightarrow \mathbb{R}_{\geq 0}$ such that for all $x,y \in S$, $d(x,y) = 0 \iff x = y$ (identity of indiscernibles) and $d(x,y) = d(y,x)$ (symmetry). Then, we call $(S, d)$ a \emph{semimetric space} and $d$ a \emph{distance} function on $S$.
\end{definition}
% lax end

Using a distance, we define local $\epsilon$-robustness at an input $x^*$ requiring that sufficiently close (similar) inputs receive the same label. If it holds for all possible inputs, the classifier is globally $\epsilon$-robust.
% lax begin Lax544225.ClassifierRobustness
\begin{definition}[$\epsilon$-Robustness]
\label{def:local-robustness}
Let $(S,d)$ be a semimetric space. A classifier $c: S \rightarrow [n]$ is \emph{locally robust} at a point $x^* \in S$ for $\epsilon > 0$ if:
$$\forall x \in S [d(x, x^*) \leq \epsilon \rightarrow c(x) = c(x^*)].$$    
\end{definition}
% lax end

If a classifier $c$ is given, we call $x \in S$ a \textit{(non-)robust point} if $c$ is (not) locally robust on $x$. If a classifier is robust at all points in $S$, it is \emph{globally} $\epsilon$-robust.

\paragraph{Research Gap.} Consider the question of defining global robustness. An obvious way is to enforce local robustness for every point. But as ref. \cite{li2024certifying} points out, \textit{"local robustness verification cannot be directly generalized to the whole decision region of a class"} because the distance can \textit{"cross decision boundaries [...], making a verification task unrealistic"}. It is easy to see that points close to the true decision boundary are not robust which makes robustness on every point infeasible. 

Most works explore global robustness under additional assumptions. Different authors restrict robustness to specific datasets or underlying data distributions~\cite{casadio2022neural,tsipras2018robustness,zhang2019theoretically,yang2020closer,mangal2023certifying}. Some authors study the trade-off between robustness and accuracy and demonstrate that, for a particular classification problem, improvements in accuracy necessarily come at the cost of robustness~\cite{tsipras2018robustness,zhang2019theoretically}.
In contrast, others show that robustness and accuracy can positively correlate when class distributions are separable, a property observed in several standard image datasets~\cite{yang2020closer,mangal2023certifying}. 

\begin{figure}[t]
\centering
\definecolor{muBlue}{HTML}{1F4E79}
\definecolor{epsOrange}{HTML}{D9730D}
\begin{tikzpicture}[scale=0.58]
  % lightly shaded l-infinity epsilon-balls of the minority inputs:
  % together they cover the whole grid
  \foreach \x in {2,7}{
    \foreach \y in {2,7}{
      \fill[orange!10] (\x-2.5,\y-2.5) rectangle (\x+2.5,\y+2.5);
    }
  }

  % faint grid to read the lattice as a grid of inputs;
  % the thick ball boundaries run between rows and columns of inputs
  \foreach \i in {0,...,9}{\draw[gray!18] (\i,-0.5)--(\i,9.5);}
  \foreach \j in {0,...,9}{\draw[gray!18] (-0.5,\j)--(9.5,\j);}

  % l-infinity epsilon-balls of the minority inputs (half-width eps=2.5)
  \foreach \x in {2,7}{
    \foreach \y in {2,7}{
      \draw[epsOrange, very thick] (\x-2.5,\y-2.5) rectangle (\x+2.5,\y+2.5);
    }
  }

  % eps arrow (half-width of an l-inf ball), drawn beneath the inputs
  \draw[<->,thick,gray!70] (2,2) -- (4.5,2);
  \node[anchor=south,inner sep=1.5pt] at (3.3,2.05) {\footnotesize $\epsilon$};

  % majority-label (class A) inputs on the whole grid
  \foreach \i in {0,...,9}{
    \foreach \j in {0,...,9}{
      \fill[muBlue] (\i,\j) circle (1.7pt);
    }
  }

  % minority-label (class B) inputs on a sparse sublattice
  \foreach \x in {2,7}{
    \foreach \y in {2,7}{
      \fill[epsOrange] (\x-0.17,\y) -- (\x,\y+0.17) -- (\x+0.17,\y) -- (\x,\y-0.17) -- cycle;
    }
  }

  % compact legend to the right of the grid
  \begin{scope}[shift={(10.7,5.7)}]
    \fill[muBlue] (0,0) circle (1.7pt);
    \node[anchor=west] at (0.35,0) {\footnotesize majority label $A$};
    \fill[epsOrange] (-0.17,-1) -- (0,-0.83) -- (0.17,-1) -- (0,-1.17) -- cycle;
    \node[anchor=west] at (0.35,-1) {\footnotesize minority label $B$};
    \draw[epsOrange, very thick, fill=orange!10] (-0.05,-2.5) rectangle (0.5,-1.95);
    \node[anchor=west] at (0.7,-2.225) {\footnotesize $\epsilon$-ball of $B$};
  \end{scope}
\end{tikzpicture}
\caption{Probabilistic robustness does not constrain the \emph{structure} of
non-robust points. In the image, inputs of class $B$ are rare, so a randomly drawn
$\epsilon$-close pair is almost always same-label and density-based
robustness bound is met. Yet the $B$ inputs sit on a sublattice whose
$\epsilon$-balls (which become squares with $\ell_\infty$ norm) cover the entire input domain: every input of
the majority class $A$ lies within $\epsilon$ of some $B$, so no input is locally robust.}
\label{fig:grid}
\end{figure}

When it comes to measuring global robustness, authors introduce additional concepts to avoid working with robustness at every point. A common approach is to enforce robustness only on points where a classifier has high confidence~\cite{athavale2024verifying,leino2021globally}. This solution creates another problem of ensuring that classifier prediction does reach high confidence whenever required~\cite{el2010foundations,chow1970optimum}. Another popular approach is to enforce that a probability of finding a pair $x, x^*\in S$ with $d(x,x^*) \leq \epsilon$ and $c(x) \neq c(x^*)$ is less than some threshold value~\cite{levy2023groma,li2024certifying,mangal2019robustness}. Probabilistic approach only bounds the density of non-robust points and does not tell anything about their structure. For example, a grid of singular points such as one in Figure~\ref{fig:grid} with a different classifier label from its neighbors can result in a satisfaction of probabilistic robustness for a classifier that has no robust points (as they have a non-robust point from the grid as a neighbor).

\paragraph{This Work: $(\mu,\epsilon)$-Robustness.} In this work, we explore a new version of global classifier robustness. We consider a classifier to be $(\mu,\epsilon)$-robust if every non-robust point has a robust point located at most $\mu$-far. Note that $(0,\epsilon)$-robustness is equivalent to $\epsilon$-robustness.
% lax begin Lax544225.ClassifierRobustness
\begin{definition}[$(\mu,\epsilon)$-Robustness]
\label{def:mueps}
Let $(S,d)$ be a semimetric space and\linebreak $c: S \rightarrow [n]$ a classifier. $c$ is $(\mu, \epsilon)$-robust for $\mu \in \mathbb{R}_{\geq0}$ and $\epsilon \in \mathbb{R}_{>0}$ if for all $x \in S$, there exists $x' \in S$ with $d(x,x') \leq \mu$ such that $c(x) = c(x')$ and $c$ is locally $\epsilon$-robust at $x'$.
\end{definition}
% lax end

This property is a novel domain-independent version of global classifier robustness. Uniquely, this property enforces additional structure on non-robust points, enforcing coverage of the input domain by robust points. In the following sections, we will explore properties of $(\mu,\epsilon)$-robust classifiers in a closer detail.
\section{Robust Partitions via Covering}
\label{sec:partitions-covering}

In this section, we pursue two goals. First, we introduce a formulation of $(\mu,\epsilon)$-robustness using partitions. We then show that $(\mu,\epsilon)$ is equivalent to covering, mirroring equivalence of $\epsilon$-robustness to separability.

\subsection{Robust Partitions}

To avoid ambiguities, we introduce some well-known definitions from the literature that are relevant to our developments.
% lax begin Lax544225.Partitions
\begin{definition}[Partition]
\label{def:partition}
A collection of nonempty sets $\{P_1,\dots,P_p\}$ is a partition of a set $S$ iff
\begin{enumerate}
    \item $\bigcup_{i \in [p]} P_i = S$ and
    \item $P_i \cap P_j = \emptyset$ for all $i \neq j$.
\end{enumerate} 
The partition $\{S\}$ on the set $S$ is called \emph{trivial}.
\end{definition}
% lax end

% lax begin Lax544225.Partitions
\begin{definition}[Classifier Induced Partition]
\label{def:induced-partition}
    A classifier $c: S \rightarrow [n]$ induces a partition $P^c := \{S^c_1, S^c_2, \dots S^c_n\}$ with $S^c_i := \{x \in S\,|\,c(x) = i\}$.
\end{definition}
% lax end

Robustness is often described as invariance under $\epsilon$-sized perturbations. This perspective allows us to reformulate robustness, for which we define the close neighborhood around an input as an $\epsilon$-ball:
% lax begin Lax544225.Balls
\begin{definition}[$\epsilon$-Ball]
\label{def:epsilon-ball}
Let $(S, d)$ be a semimetric space. Then $B(x^*, \epsilon) := \{x \in S\,|\,d(x,x^*) \leq \epsilon\}$ is an \emph{$\epsilon$-ball} around $x^*$.
\end{definition}
% lax end

Despite the name, $B(x^*, \epsilon)$ may not resemble a geometric ball. In fairness notions, for example, the distance function may capture similarity in qualification of individuals rather than geometric proximity of their attribute vectors.
% lax begin Lax544225.PartitionRobustness
\begin{definition}[$(\mu,\epsilon)$-Robustness of a Partition]
\label{def:mueps-robustness-partition}
Let $(S,d)$ be a semimetric space. Then a partition $P := \{S_1,\dots,S_p\}$ of $S$
is \emph{globally robust} for $\mu \geq 0, \epsilon > 0$ if and only if:
$$\forall x \in S\,\exists x' \in B(x,\mu) \,[\bigvee_{i \in [p]} (\{x\} \cup B(x', \epsilon) \subseteq S_i)].$$
\end{definition}
% lax end

This subtle reformulation completely shifts $(\mu,\epsilon)$-robustness to a property of the induced partition. 
By working with partitions, our framework does not require the number of classes to be specified a priori, unlike some existing approaches~\cite{yang2024geometric}. To streamline subsequent results, we will use the term \emph{$(\mu,\epsilon)$-robust partition} for partitions that satisfy Definition~\ref{def:mueps-robustness-partition}. If parameters are clear from the context, we will omit $(\mu,\epsilon)$.

\subsection{From Separability to Covering}

A folklore fact about global robustness is its equivalence to \emph{$\epsilon$-separability}~\cite{yang2020closer,mangal2023certifying} and we can to extend this property to our new robustness notion, albeit in a more complicated form. 
% lax begin Lax544225.CoveringCharacterization
\begin{theorem}
\label{thm:separability}
    Let $(S,d)$ a semimetric space, $P := \{S_1, \dots, S_n\}$ a partition of $S$ and $\mu \geq 0, \epsilon > 0$. Then the following are equivalent:
% lax end
    \begin{enumerate}
        \item $P$ is $(\mu,\epsilon)$-robust.
% lax begin Lax544225.Covering
        \item for all $S_i \in P$ and all $x \in S_i$ holds that $B(x, \mu) \not\subseteq \bigcup_{x^* \in S \setminus S_i} B(x^*, \epsilon)$.
% lax end

    \end{enumerate}
\end{theorem}
\begin{proof}
    We prove both directions separately.

    $(1.) \Rightarrow (2.)$: 
    For contradiction, assume that (1.) holds but (2.) does not. In this case, there exists $S_i \in P$ with $x \in S_i$ such that $B(x,\mu) \subseteq \bigcup_{x^* \in S \setminus S_i} B(x^*, \epsilon)$. From Defintion~\ref{def:mueps-robustness-partition}, there exists $x' \in B(x,\mu)$ with $B(x',\epsilon) \subseteq S_i$. From symmetry of distance (Definition~\ref{def:distance}), it follows that $x' \notin \bigcup_{x^* \in S \setminus S_i} B(x^*, \epsilon)$, leading to a contradiction with (1.).

    $(2.) \Rightarrow (1.)$:
    Similarly, assume that (2.) holds but not (1.). By Defintion~\ref{def:mueps-robustness-partition}, there exists $x \in S$ such that for all $x' \in B(x,\mu)$ have $B(x', \epsilon) \not\subseteq S_i$ for all $S_i \in P$. Pick a specific $S_i' \in P$ such that $x \in S_i'$. By symmetry of distance, each $x' \in B(x,\mu)$ is also contained in some $B(x^*,\epsilon)$ with $x^* \notin S_i'$. This implies $B(x, \mu) \subseteq \bigcup_{x^* \in S \setminus S_i'} B(x^*, \epsilon)$, raising a contradiction to (2.).
    \qed
\end{proof}

\begin{figure}[t]
  \centering
  \definecolor{muBlue}{HTML}{1F4E79}
  \definecolor{epsGreen}{HTML}{2E8B57}
  \definecolor{epsOrange}{HTML}{D9730D}
  \begin{tikzpicture}[scale=1.5, every node/.style={transform shape},
      classLine/.style = {black!75, line width=0.9pt},
      ballS2/.style    = {draw=epsGreen,  thick},
      ballS3/.style    = {draw=epsOrange, thick},
      dotS1/.style     = {circle, fill=muBlue,    inner sep=1.5pt},
      dotS2/.style     = {circle, fill=epsGreen,  inner sep=1.5pt},
      dotS3/.style     = {circle, fill=epsOrange, inner sep=1.5pt},
      radArrow/.style  = {-{Stealth[length=2.2mm]}, thin, gray!70},
    ]
    % ----------------------------------------------------------------
    % Geometry (numerically verified):
    %   decision boundary (vertical, wiggly):
    %     x(y) = 0.55 - 0.85*cos(3y)   [3 rad = 171.9 degrees]
    %   S_1 = right of the curve, S_2/S_3 = left of it (split at y=-0.1)
    %   x   = (0, 0) in S_1,  mu  = 1.1
    %   eps = 1.0   (mu = 1.1*eps > eps)
    %   eps-balls (radius 1.0), centers in S_2 / S_3:
    %     x_1 = ( 0.28,  0.48)  in S_2
    %     x_2 = (-0.55,  0.00)  in S_2
    %     x_3 = ( 0.28, -0.48)  in S_3
    %   The three eps-balls cover the mu-ball of x completely
    %   (max distance of any mu-ball point to the nearest center: 0.953);
    %   neither S_2 alone nor S_3 alone covers it.
    %   Note: three eps-disks can cover a mu-disk of radius at most
    %   2*eps/sqrt(3) ~ 1.155*eps, so mu = 1.1*eps is near the
    %   largest ratio possible with three covering balls.
    %   Canvas is cropped tight: x in [-1.75,1.5], y in [-1.7,1.7].
    % ----------------------------------------------------------------
    \coordinate (x)  at (0,0);
    \coordinate (x1) at (0.28,0.48);
    \coordinate (x2) at (-0.55,0);
    \coordinate (x3) at (0.28,-0.48);

    % ---- landscape: the three class regions ---------------------------
    \fill[blue!8]                                            % S_1 (right)
      plot[domain=1.7:-1.7, samples=250, smooth]
        ({0.55-0.85*cos(171.9*\x)},{\x})
      -- (1.5,-1.7) -- (1.5,1.7) -- cycle;
    \fill[green!12]                                          % S_2 (upper left)
      (-1.75,1.7)
      -- plot[domain=1.7:-0.1, samples=150, smooth]
           ({0.55-0.85*cos(171.9*\x)},{\x})
      -- (-1.75,-0.1) -- cycle;
    \fill[orange!10]                                         % S_3 (lower left)
      (-1.75,-0.1)
      -- plot[domain=-0.1:-1.7, samples=150, smooth]
           ({0.55-0.85*cos(171.9*\x)},{\x})
      -- (-1.75,-1.7) -- cycle;

    % ---- coverage blankets: union of the eps-balls per class ----------
    % (transparency group per class => each class gets ONE uniform tint)
    \begin{scope}[transparency group, opacity=0.2]
      \fill[epsGreen] (x1) circle (1.0);
      \fill[epsGreen] (x2) circle (1.0);
    \end{scope}
    \begin{scope}[transparency group, opacity=0.2]
      \fill[epsOrange] (x3) circle (1.0);
    \end{scope}

    % ---- class boundaries ----------------------------------------------
    \draw[classLine]
      plot[domain=-1.7:1.7, samples=250, smooth]
        ({0.55-0.85*cos(171.9*\x)},{\x});
    \draw[classLine] (-1.75,-0.1) -- (-0.26,-0.1);

    % ---- highlighted epsilon-balls ----------------------------------------
    \draw[ballS2] (x1) circle (1.0);
    \draw[ballS2] (x2) circle (1.0);
    \draw[ballS3] (x3) circle (1.0);

    % ---- mu-ball around x (mu > eps) --------------------------------------
    \draw[muBlue, very thick] (x) circle (1.1);

    % ---- points --------------------------------------------------------------
    \node[dotS1] at (x)  {};
    \node[dotS2] at (x1) {};
    \node[dotS2] at (x2) {};
    \node[dotS3] at (x3) {};

    % ---- radius annotations -----------------------------------------------------
    \draw[radArrow] (x)  -- ($(x)+(250:1.1)$)  node[midway, left]        {$\mu$};
    \draw[radArrow] (x1) -- ($(x1)+(60:1.0)$)  node[midway, above right] {$\epsilon$};
    \draw[radArrow] (x2) -- ($(x2)+(150:1.0)$) node[midway, above left] {$\epsilon$};
    \draw[radArrow] (x3) -- ($(x3)+(300:1.0)$) node[midway, below right] {$\epsilon$};

    % ---- point labels ----------------------------------------------------------------
    \node[below right=0pt]  at (x)  {$x$};
    \node[above right=1pt]  at (x1) {$x_1$};
    \node[left=2pt]         at (x2) {$x_2$};
    \node[below right=1pt]  at (x3) {$x_3$};

    % ---- region labels --------------------------------------------------------------------
    \node[font=\small] at (-1.5, 1.45) {$S_2$};  % upper left
    \node[font=\small] at (-1.5,-1.3)  {$S_3$};  % lower left
    \node[font=\small] at ( 1.25, 1.45) {$S_1$}; % right
  \end{tikzpicture}
  \caption{Illustrated violation of $(\mu,\epsilon)$-robustness though covering condition (2.) of Theorem~\ref{thm:separability}: the $\mu$-ball of $x \in S_1$ is fully covered by three $\epsilon$-balls of $x_1, x_2 \in S_2$ and $x_3 \in S_3$.}
  \label{fig:covering}
\end{figure}

The condition (2.) in Theorem~\ref{thm:separability} can be stated as: no $\mu$-neighbourhood is covered by $\epsilon$-neighbourhoods of points with a different label. An example when it is not the case is illustated in Figure~\ref{fig:covering}. As $(\mu, \epsilon)$-robustness is a generalization of $\epsilon$-robustness, our proof is also a formalization of separability for $\epsilon$-robustness. For $\mu = 0$, the condition (2.) reduces to each point (0-neighbourhood) being $\epsilon$-far from any point with a different label. Separability is established here in a general, non-probabilistic setting without distributional assumptions. Unlike prior work, we show that separability is a \emph{necessary} condition for global robustness. 

The characterization of robustness from Theorem~\ref{thm:separability} is central to future results and we will use it as a primary criterion of partition robustness.


\section{Modeling Finite Precision Arithmetic}

To bridge theory and practice, we want to formulate our results on a structure that exists under computer arithmetic constraints. Modeling classifiers over real numbers can lead to a mismatch between theoretical insights and practical implementations. For example, verifiers for neural networks based on real arithmetic can be unsound when applied to floating-point computations~\cite{jia2021exploiting,cordeiro2025neural}.

\paragraph{Mathematcal Model.} Floating point numbers approximate real values with finite precision according to the IEEE 754 standard, which defines a finite arithmetic over a set of bitvector-encoded numbers.
Reasoning about finite-precision arithmetic with standard operators (such as addition) is challenging, since these operators may violate familiar algebraic properties like associativity~\cite{Kneusel2025}. 
An additional problem is presented by special floating point values such as \texttt{NaN} or \texttt{Inf}.
Nevertheless, the IEEE 754-2008 standard provides a well-defined total ordering on floating-point numbers~\cite{4610935}, implemented, for example, in Rust\footnote{\url{https://doc.rust-lang.org/std/primitive.f64.html\#method.total\_cmp}}. This motivates formulation of our results over ordered spaces $(S, <, d)$.

But an introduction of an order will not lead to more productive results unless a relationship between the order and the distance exists. To fix this, we introduce a new inequality. The inequality states that visiting an intermediate point between two points in the order produces at most the direct distance. In some sense, it is the opposite of a triangle inequality, which states that visiting an additional point increases the distance. But as our goal in further proofs is to bound the size of balls, we found this inequality more productive. We also decided against making this an equality for a irrelevant to our further results possibility that distance uses floating point rounding.

With these considerations, we arrive at our mathematical model standing for computer arithmetic spaces:
% lax begin Lax544225.OrderedSemimetric
\begin{definition}[Ordered Semimetric Space]
    \label{def:ordered-space}
    A tuple $(S, <, d)$ with a set $S$, a (partial) order $< \subseteq S \times S$, and $d: S \times S \rightarrow \mathbb{R}_{\geq 0}$ is an \emph{ordered semimetric space} if $(S, d)$ is a semimetric space and for all $a,b,c \in S$ such that $a \leq b \leq c$ implies $d(a, c) \geq d(a, b) + d(b,c)$.
\end{definition}
% lax end


\paragraph{Model Validation.} After defining ordered semimetric spaces, the question remains whether such spaces indeed exist on floating point numbers and can resemble common distances. To verify this, we utilized Mistral Vibe to generate provably correct instances of ordered semimetric spaces with respect to the IEEE 754 total order on lists of Flocq~\cite{boldo2011flocq} floating point numbers within the Rocq proof asisstant~\cite{rocq}, re-implement the resulting distances for Python  64 bit \texttt{float} and produce a script for differential testing between Rocq and Python using fixed edge cases and randomly generated instances. The task was successfully done by Mistral Vibe in 6 hours and did cost approximately 45 euro. The output was reviewed by humans afterwards and is avaliable on GitHub\footnote{\url{https://github.com/verinncoq/mueps-robustness-proofs/tree/main/flocq_ordered_spaces}}. The following proposition summarizes the findings.

\begin{proposition} 
\label{prop:fp-space}
% lax begin Lax544225.BinaryFloats
Let $\mathcal{F}$ be a binary floating point format with $p \in \mathbb{N}$ being the number of mantissa bits. Then, it consists of triples $(s, c, m)$ representing numbers $(-1)^s \cdot 2^c \cdot m$ such that $s \in \{0,1\}$, $c \in \{c_{\min},\dots, c_{\max}\} \subset \mathbb{Z}$, $m \in \{k \cdot 2^{-p} : k \in \{1, \dots, 2^{p+1}-1\} \subset \mathbb{Z}\}$ together with additional values $\pm0$, \texttt{NaN}, $\pm$\texttt{Inf} $\in \mathcal{F}$. For each $(s, c, m) \in \mathcal{F}$ holds that either $m \geq 1$ (normal) or $c = c_{\min}$ (subnormal). Let $\phi : \mathcal{F} \rightarrow \mathbb{R}$ with:
\begin{itemize}
    \item $\phi((s,c,m)) = (-1)^s \cdot 2^c \cdot m$ 
    \item $\phi(+0) = 0$ and $\phi(-0) = -2^{c_{\min} - p - 1}$
    \item $\phi(\texttt{NaN}) = 2^{c_{\max} + 1}  + 1$
    \item $\phi(\pm\texttt{Inf}) = \pm 2^{c_{\max} + 1}$
\end{itemize}
% lax end
Then, the following are ordered semimetric spaces:
\begin{enumerate}
% lax begin Lax544225.FloatAbsoluteSpace
    \item $(\mathcal{F}, <_{IEEE}, d_1)$ where $<_{IEEE}$ is the IEEE total order and $d_1(f_1, f_2) = |\phi(f_1) - \phi(f_2)|$.
% lax end
% lax begin Lax544225.FloatSquaredSpace
    \item $(\mathcal{F}, <_{IEEE}, d_2)$ where $d_2(f_1, f_2) = (\phi(f_1) - \phi(f_2))^2$.
% lax end
% lax begin Lax544225.ProductSpaces
    \item $(\mathcal{F}^n, <^n, d_3(f, g) := \sum_{i \in [n]} d(f_i, g_i))$
% lax end
% lax begin Lax544225.SumProductSpace
 such that $(\mathcal{F}, <, d)$ is an ordered semimetric space and $<^n$ is n-wise product partial order.
% lax end
% lax begin Lax544225.MaximumProductSpace
    \item $(\mathcal{F}^n, <^{n}_*, d_4(f, g) := \max_{i \in [n]}d(f_i, g_i))$ where $(\mathcal{F},<,d)$ is an ordered semimetric space such that for all $a,b,c\in\mathcal{F}$ with $a \leq b \leq c$ holds that $d(a,c) = d(a,b) + d(b,c)$ and $<^n_* := \{(f,g) \in \mathcal{F}^n \times \mathcal{F}^n\,|f <^n g \wedge \forall i,j \in [n]\, d(f_i, g_i) = d(f_j, g_j)\,\}$.
% lax end
    \item Extensions of
% lax begin Lax544225.SumListSpace
(3)
% lax end
and
% lax begin Lax544225.MaximumListSpace
(4)
% lax end
% lax begin Lax544225.ListSpaces
to arbitrary finite lists of floating-point numbers where a shorter list is padded with \texttt{NaN} for distance computation, and a penalty equal to $1$ is added to the distance whenever list lengths are different. The partial orders compare only lists of the same length, and lists of different lengths are incomparable.
% lax end
\end{enumerate}
\end{proposition}

The combinations of spaces 1 and 3, 1 and 4 in Proposition~\ref{prop:fp-space} give us notions similar to $l_1$, $l_\infty$ distances respectively. Overall, we showed that ordered semimetric spaces can be faithfully constructed on floating point numbers, showing that the results proven in the following sections are transferrable to computer arithmetic.
\section{Existence of Robust Partitions}
\label{sec:existence}

The main criticism of $\epsilon$-robustness is that it is impossible to satisfy as separability cannot be enforced on a decision boundary. We will explore our two-fold answer to this critique. First, we should notice that the non-satisfiability argument only applies to continuous domains such as real numbers while most classifiers operate on a discrete domain of floating-point numbers. Second, this problem gets alleviated for $(\mu,\epsilon)$-robustness. To support two arguments with mathematical evidence, we will prove conditions for existence of at least one non-constant robust classifier partition. In Section~\ref{subsec:general-existence}, we derive a general criterion on an arbitrary semimetric space and in Section~\ref{subsec:existence-order}, we specialize it to ordered semimetric spaces, thus creating a connection to finite precision arithmetic.

\subsection{General Existence Criterion}
\label{subsec:general-existence}

We study the conditions under which $(\mu, \epsilon)$-robustness is satisfiable by a non-constant classifier, i.e., one that distinguishes between at least two classes. The following result shows that $(\mu, \epsilon)$-robustness imposes structural constraints on the domain.
% lax begin Lax544225.GeneralExistence
\begin{theorem}[General Existence]
\label{thm:existence}
Let $(S,d)$ a semimetric space and $\mu \geq 0, \epsilon > 0$. The follwing statements are equivalent:
\begin{enumerate}
    \item The only $(\mu,\epsilon)-$robust partition on $S$ is trivial.
    \item For all nonempty $S' \subset S$, there exists $x \in S$ such that either $x \in S'$ and $B(x,\mu) \subseteq \bigcup_{x' \in S \setminus S'} B(x', \epsilon)$ or $x \in S \setminus S'$ with $B(x,\mu) \subseteq \bigcup_{x' \in S'} B(x', \epsilon)$.
\end{enumerate}
\end{theorem}
% lax end

\begin{proof}
    We prove each direction separately:

    $(1.) \Rightarrow (2.)$:
    Assume (1.) and let $S' \subset S$ be a proper subset such that condition (2.) is false. We prove that $P := \{S', S\setminus S'\}$ is $(\mu, \epsilon)$-robust, leading to contradiction with (1.). If $P$ is not robust, then by Theorem~\ref{thm:separability} there exists $x \in S$ with either $x \in S'$ and $B(x, \mu) \subseteq \bigcup_{x' \in S \setminus S'} B(x', \epsilon)$ or $x \in S \setminus S'$ with $B(x,\mu) \subseteq \bigcup_{x' \in S'} B(x', \epsilon)$. Note that it is exactly the condition (2.) of this theorem, which was assumed to be false, implying by contradiction that $P$ is a robust partition.

    $(2.) \Rightarrow (1.)$:
    Assume (2.) and let $P := \{S_1, \dots, S_n\}$ with $n > 1$ be any robust partition of $S$. For contradiction, we show that $S_1$ violates (2.). By condition (2.) from Theorem~\ref{thm:separability}, there is no $x \in S_1$ with $B(x, \mu) \subseteq \bigcup_{x' \in S \setminus S_1} B(x', \epsilon)$. It leaves the case that $x \in S \setminus S_1$ with $B(x,\mu) \subseteq \bigcup_{x' \in S_1} B(x', \epsilon)$. This means that $x \in S_i$ for some $ i \in \{2,\dots,n\} \subset [n]$ and by Theorem~\ref{thm:separability}, $B(x, \mu) \not\subseteq \bigcup_{x' \in S\setminus S_i} B(x', \epsilon)$. As $\bigcup_{x' \in S_1} B(x', \epsilon) \subseteq \bigcup_{x' \in S\setminus S_i} B(x', \epsilon)$, $S_1$ violates condition (2.) of this theorem in all cases.
    \qed
\end{proof}

Given fixed $\mu$ and $\epsilon$, Theorem~\ref{thm:existence} establishes a class of semimetric spaces that have at least one nontrivial robust partition. Contrary to the common perception that robustness is impossible for $\mu = 0$, many sets and distance functions violate condition (2) of Theorem~\ref{thm:existence}, e. g., consider the set $\{0,2,4,6\}$ with the standard metric $d(x,y) = |x - y|$ and let $\epsilon = 1$. This observation demonstrates benefits of providing mathematical clarity by replacing informal arguments with rigorous analysis.

\subsection{Existence of a Robust Partition on a Total Order}
\label{subsec:existence-order}

% lax begin Lax544225.Boundary
Denote by $Succ(x,y)$ a formula that is true if $x < y$ and there is no $z$ with $x < z < y$. We start with results on finite totally ordered semimetric spaces $(S, <, d)$ where $S$ is arbitrary and the order $<$ is total. For each subset $S' \subseteq S$ of such a space, a \emph{boundary} $\partial S' = \{x \in S' \,|\, \exists y \not\in S' . Succ(x,y) \vee Succ(y,x)\}$ and \emph{complement boundary} $\overline\partial S' = \{x \not\in S' \,|\, \exists y \in S' . Succ(x,y) \vee Succ(y,x)\}$.
% lax end


% lax begin Lax544225.BoundaryCoveringCharacterization
\begin{theorem}
\label{thm:existence-total}
    Let $(S, <, d)$ be a finite totally ordered semimetric space and $\mu \geq 0, \epsilon > 0$. The following two statements are equivalent:
    \begin{enumerate}
        \item For all nonempty $S' \subset S$, there exists $x \in S$ such that either $x \in S'$ and $B(x,\mu) \subseteq \bigcup_{x' \in S \setminus S'} B(x', \epsilon)$ or $x \in S \setminus S'$ with $B(x,\mu) \subseteq \bigcup_{x' \in S'} B(x', \epsilon)$.
        \item For all nonempty $S' \subset S$, there exists $x \in S$ such that either $x \in S'$ and $(B(x,\mu) \cap S') \subseteq \bigcup_{x' \in \overline{\partial}S'} B(x', \epsilon)$ or $x \in S \setminus S'$ with $(B(x,\mu) \cap (S \setminus S')) \subseteq \bigcup_{x' \in \partial S'} B(x', \epsilon)$.
    \end{enumerate}
\end{theorem}
% lax end

\begin{proof}
    Again, we prove each direction separately:

    \textbf{(1.) $\Rightarrow$ (2.)}:
    To prove (2.) of this theorem, pick any nonempty $S' \subset S$. By applying (1.) to $S'$, we get two cases:
    \begin{itemize}    
        \item \textit{1. Case: $x \in S'$ and $B(x,\mu) \subseteq \bigcup_{x' \in S \setminus S'} B(x', \epsilon)$}. To show (2.), we will prove that $(B(x,\mu) \cap S') \subseteq \bigcup_{x' \in \overline{\partial}S'} B(x', \epsilon)$. Pick some $x^* \in B(x, \mu) \cap S'$ and select some $x'' \in S \setminus S'$ with $x^* \in B(x'', \epsilon)$, which exists by assumption. W.l.o.g. assume that $x^* < x''$, implying existence of a $x_{\overline{\partial}} \in \overline{\partial} S'$ with $x^* \leq x_{\overline{\partial}} \leq x''$. Then, $d(x^*, x_{\overline{\partial}}) \leq d(x^*, x'') \leq \epsilon$ by Definition~\ref{def:ordered-space}. From this, $x^* \in B(x_{\overline{\partial}}, \epsilon)$ follows. As such $x_{\overline{\partial}} \in \overline{\partial} S'$ exists for each $x^* \in B(x, \mu) \cap S'$, we deduce $(B(x,\mu) \cap S') \subseteq \bigcup_{x' \in \overline{\partial}S'} B(x', \epsilon)$, and (2.) is proven.
        \item \textit{2. Case: $x \in S \setminus S'$ and $B(x,\mu) \subseteq \bigcup_{x' \in S'} B(x', \epsilon)$.} Proven analogously to the 1. Case by finding $x_{\partial} \in \partial S'$ for each $x^* \in B(x,\mu) \cap S \setminus S'$ for some $x \in S$.
    \end{itemize}

    \textbf{(2.) $\Rightarrow$ (1.)}:
    To prove (1.), pick a nonempty $S' \subset S$. By (2.), there exists $x \in S$ satisfying one of two conditions in (2.). We get two cases:
    \begin{itemize}
        \item \textit{1. Case: $x \in S'$ and $(B(x, \mu) \cap S') \subseteq \bigcup_{x' \in \overline{\partial}S'}  B(x', \epsilon)$}. Consider some $x^* \in B(x, \mu)$. If $x^* \in S'$ then by (2.), $x^* \in \bigcup_{x' \in \overline{\partial}S'} B(x', \epsilon) \subseteq \bigcup_{x' \in S \setminus S'} B(x', \epsilon)$. If $x^* \in S \setminus S'$ then $x^* \in B(x^*, \epsilon) \subseteq \bigcup_{x' \in S \setminus S'} B(x', \epsilon)$. Combined, we deduce $B(x, \mu) \subseteq \bigcup_{x' \in S \setminus S'} B(x', \epsilon)$ and (1.) is proven.
        \item \textit{2. Case: $x \in S \setminus S'$ and $(B(x,\mu) \cap (S \setminus S')) \subseteq \bigcup_{x' \in \partial S'} B(x', \epsilon)$.} Proven similarly to the 1. Case.
    \end{itemize}
    \qed
\end{proof}

% lax begin Lax544225.BoundaryExistence
\begin{corollary}
    \label{cor:existence-total}
    Let $(S, <, d)$ a finite totally ordered semimetric space and $\mu \geq 0, \epsilon > 0$. The following two statements are equivalent:
    \begin{enumerate}
        \item The only robust partition on $S$ is trivial.
        \item For all nonempty $S' \subset S$, there exists $x \in S$ such that either $x \in S'$ and $(B(x,\mu) \cap S') \subseteq \bigcup_{x' \in \overline{\partial}S'} B(x', \epsilon)$ or $x \in S \setminus S'$ with $(B(x,\mu) \cap (S \setminus S')) \subseteq \bigcup_{x' \in \partial S'} B(x', \epsilon)$.
    \end{enumerate}
\end{corollary}
% lax end

\begin{proof}
    Note that condition (1.) in Theorem~\ref{thm:existence-total} is exactly condition (2.) in Theorem~\ref{thm:existence}.
\end{proof}

With this, we eastablished necessary and sufficient condition for existence of a nontrivial robust partition in an ordered semimetric space. Note that as Theorem~\ref{thm:existence-total} is an equivalence condition, the condition (2.) cannot be simplified further. A covered $B(x, \mu)$ is not necessarily produced by $x \in \partial S'$ or $x \in \overline{\partial} S'$. Also, multiple $\epsilon$-balls may be required to cover a $\mu$-ball. In summary, we can observe that studying equivalent conditions to existence does not lead to simple conditions. Due to this observation, we will switch to estimating upper bounds on parameters $\mu$ and $\epsilon$ for nontrivial robust partitions.




\section{Precision-Related Upper Bounds}
\label{sec:upper-bounds}

An interesting empirically supported fact for $\epsilon$-robustness is that quantization, a process of reducing numerical precision of classifier's inputs and parameters, causes classifier to become more robust~\cite{xiao2023robustmq,shrestha2024properties,duncan2020relative}. The explanation is simple: larger distances between input space points raises the threshold on admissible $\epsilon$-values and thereby makes $\epsilon$-robustness easier to satisfy. Mathematically, one can show that with larger gaps maximum $\epsilon$ for nontrivial robust partitions increases. In this section, our goal is to derive such upper bounds on $\mu$ and $\epsilon$ based on precision measures: in Section~\ref{subsec:bounds-total} for totally ordered semimetric spaces and in Section~\ref{subsec:bounds-partial} for general partial orders such as multidimensional inputs common to classifiers such as neural networks.

\subsection{Upper Bounds on Totally Ordered Spaces}
\label{subsec:bounds-total}

We start with upper bounds on totally ordered semimetric spaces introduced in Section~\ref{subsec:existence-order}. To state our results, we define characteristics of distances between two distinct representable numbers in an order representing finite-precision arithmetic. For a totally ordered semimetric space $(S,<,d)$, $a, b \in S$ with $a \leq b$, define an \textit{interval} $[a, b] := \{a \leq x \leq b\,|\,x \in S\}$.
% lax begin Lax544225.PrecisionGaps
\begin{definition}[Precision Gap Metrics]
\label{def:precision_gap}
    Let $(S, <, d)$ be a finite totally ordered semimetric space with $|S| \geq 2$. We define the metrics $\Delta$ and $\delta$ as:
    \begin{align*}
        \Delta &:= \max(\{\frac{d(x,y)}{|[x,y]| -1}\,|\,x, y \in S \wedge x < y\}) \\
        \delta &:= \max(\{\frac{|[x,y]| - 1}{d(x,y)}\,|\,x, y \in S \wedge x < y\})
    \end{align*}
    where $|\cdot|$ is set cardinality.
\end{definition}
% lax end

For floating-point numbers when $d(x,y) = |x - y|$ is used as the distance, $\Delta$ corresponds to the maximal and $\delta$ to the reciprocal of minimal \emph{unit in the last place} (ULP)~\cite{deDinechin2024}. Now, we show how $\Delta$ and $\delta$ define maximum possible $\mu$ and $\epsilon$ for nontrivial robust partitions.
% lax begin Lax544225.SuccessorCovering
\begin{theorem}
\label{thm:bounds-total}
    Let $(S, <, d)$ a finite totally ordered semimetric space and $\mu \geq 0, \epsilon > 0$. If $\epsilon \geq \Delta \cdot \delta \cdot \mu + \Delta$ then for all $x, x' \in S$ with $Succ(x,x')$, it holds that $B(x, \mu) \subseteq B(x', \epsilon)$. 
\end{theorem}
% lax end

\begin{proof}
    First, observe that $B(x, \mu) \subseteq B(x', \epsilon)$ can be restated as: for all $x^* \in S$ with $d(x^*,x) \leq \mu$ holds $d(x^*,x') \leq \epsilon$. Now, fix any $x^*, x, x' \in S$ with $Succ(x,x')$. If $x^* = x$ then $d(x^*, x') = d(x, x') \leq \Delta \leq \epsilon$. If $x^* = x'$, $d(x^*, x') = 0 \leq \epsilon$. We are left with two cases:
    \begin{itemize}
        \item \textit{1. Case: $x^* < x < x'$}: 
                    \begin{align*}
                        d(x^*,x') = & \frac{d(x^*,x')}{|[x^*,x']| - 1} \cdot (|[x^*, x']| - 1) \leq \Delta \cdot (|[x^*, x']| - 1) = \\
                                 & \Delta \cdot |[x^*,x]| = \Delta \cdot (|[x^*,x]| - 1) + \Delta = \\
                                 & \Delta \cdot \frac{|[x^*,x]| - 1}{d(x^*,x)} \cdot d(x^*,x) + \Delta \leq \Delta \cdot \delta \cdot d(x^*,x) + \Delta \leq \\
                                 & \Delta \cdot \delta \cdot \mu + \Delta \leq \epsilon
                    \end{align*}
        \item \textit{2. Case: $x < x' < x^*$}:
                    \begin{align*}
                        d(x^*, x') = &d(x', x^*) \leq d(x, x^*) \\
                                    & \leq \mu = \frac{d(x,x')}{d(x,x')} \cdot \mu \leq \Delta \cdot \delta \cdot \mu < \epsilon
                    \end{align*}
    \end{itemize}
    \qed
\end{proof}
% lax begin Lax544225.PrecisionObstruction
\begin{corollary}
\label{cor:bounds-total}
    Let $(S, <, d)$ a finite totally ordered semimetric space and $\mu \geq 0, \epsilon > 0$. If $\epsilon \geq \Delta \cdot \delta \cdot \mu + \Delta$ then the only $(\mu, \epsilon)$-robust partition is trivial. 
\end{corollary}
% lax end

\begin{proof}
   We show that conclusion of the Theorem~\ref{thm:bounds-total} implies condition (2.) of Theorem~\ref{thm:existence-total}. 
   Fix a nonempty $S' \subset S$. 
   Let $x_{min} := min(S \setminus S')$ and $x'_{min} := min(S')$.
   Since neither $S \setminus S'$ nor $S'$ are empty, both $x_{min}$ and $x'_{min}$ exist.
   \begin{itemize}
       \item \textit{1. Case: $x_{min} < x'_{min}$.}
                   Let $x_{pre} := max(\{x \in S \setminus S'\,|\, x < x'_{min}\})$ and $x'_{succ} := x'_{min}$.
                   As $Succ(x_{pre}, x'_{succ})$ holds, $B(x_{pre}, \mu) \subseteq B(x'_{succ}, \epsilon)$ by Theorem~\ref{thm:bounds-total}. 
                   Observe that $x'_{succ} \in \partial S'$. 
                   Condition (2.) of Theorem~\ref{thm:existence-total} holds as $x_{pre} \in S \setminus S'$ and $(B(x_{pre}, \mu) \cap S \setminus S') \subseteq B(x_{pre}, \mu) \subseteq B(x'_{succ}, \epsilon) \subseteq \bigcup_{x' \in \partial S'} B(x', \epsilon)$.
        \item \textit{2. Case: $x_{min} > x'_{min}$.}
                   Let $x'_{pre} := max(\{x' \in S'| x' < x_{min}\})$ and $x_{succ} := x_{min}$, then $Succ(x'_{pre}, x_{succ})$. 
                   Similarly to the 1. Case, $(B(x'_{pre}, \mu) \cap S') \subseteq B(x'_{pre}, \mu) \subseteq B(x_{succ}, \epsilon) \subseteq \bigcup_{x' \in \overline{\partial} S'} B(x', \epsilon)$.
   \end{itemize}
   Corollary~\ref{cor:existence-total} states that condition (2.) is equivalent to $(S,<,d)$ having the trivial partition as the only $(\mu, \epsilon)$-robust partition.
\end{proof}

For $\epsilon$-robustness ($\mu=0$), we see that $\epsilon \geq \Delta$ is an upper bound for a nontrivial robust partition. With increasing $\Delta$, this bound also increases which provides theoretical evidence for better robustness after quantization. Introduction of a new parameter $\mu$ allows to relax strong restrictions imposed by $\epsilon$ on partitions, which is the goal of having a weaker version of robustness. 

\subsection{Upper Bounds on Partial Orders}
\label{subsec:bounds-partial}

Usually, classifiers are not defined on a single floating-point number line but on multi-dimensional vectors that potentially consist of numbers with different precision (e.g. integers and floats). The total order given by an arithmetic over each vector dimension creates a partial order over the entire vector space. This multi-dimensional setting demands an extended result on the achievable robustness of a classifier. Unfortunately, $(\mu,\epsilon)$-robustness shows non-compositional behavior which we want to illustrate with a lemma:
% lax begin Lax544225.Noncompositionality
\begin{lemma}
\label{lem:noncompositional}
Given a semimetric space $(S, d)$ let $S_1, S_2 \subset S$ with $S_1 \cap S_2 \neq \emptyset$ as well as $\mu \geq 0, \epsilon > 0$. The following statements are \textbf{false}:
\begin{enumerate}
    \item If the only $(\mu, \epsilon)$-robust partition on $S_1$ is trivial, all $(\mu, \epsilon)$-robust partitions $P$ on $S$ include a superset of $S_1$.

    \item If the only $(\mu, \epsilon)$-robust partitions on $S_1$ and $S_2$ trivial, the only $(\mu, \epsilon)$-robust partition on $S_1 \cup S_2$ is trivial.

\end{enumerate}
\end{lemma}
% lax end
\begin{proof}
We provide counterexamples for $S := \{a,b,c,d\}$ with $d(a,b) = d(b,c) = d(c,d) = 1$, $d(a,c) = d(b,d) = 2$ and $d(a,d) = 3$. Note that $S$ has a nontrivial robust partition for $\mu = \epsilon = 1$: $\{\{a,b\},\{c,d\}\}$.
\begin{enumerate}
    \item $S_1 = \{b,c\}$ does not have a nontrivial robust partition for $\mu = \epsilon = 1$, and the nontrivial robust partition on $S$ does not contain a superset of $S_1$.
    \item $S_1 = \{a,b,c\}$ and $S_2 = \{b,c,d\}$ do not have a nontrivial robust partition for $\mu = \epsilon = 1$, but $S$ does.
\end{enumerate}
\qed
\end{proof}

Still, it is possible to find a general but relatively weak decomposition principle. First, we define how the decomposing structure looks like:
% lax begin Lax544225.Decomposition
\begin{definition}[$(\mu, \epsilon)$-Decomposition]
    \label{def:decomposition}
    A collection of nonempty sets $D := \{D_1,\dots,D_p\}$ is a $(\mu, \epsilon)$-decomposition of a semimetric space $(S, d)$ iff
    \begin{enumerate}
        \item for each $x \in D_i$, $B(x, \mu) \subseteq D_i$ and
        \item for each $x_1, x_2 \in S$, there is a sequence of sets $D_{i_1},\dots,D_{i_n} \in D$ such that $D_{i_{j}} \cap D_{i_{j+1}} \neq \emptyset$ and $x_1 \in D_{i_1}$, $x_2 \in D_{i_n}$. 
    \end{enumerate} 
\end{definition}
% lax end


Now, we are ready to show that $(\mu, \epsilon)$-decompositions allow us to focus on considering $(\mu, \epsilon)$-robustness of individual decomposition elements only.

% lax begin Lax544225.DecompositionObstruction
\begin{theorem}
\label{thm:decomposition}
    Given a semimetric space $(S, d)$ with $\mu \geq 0, \epsilon > 0$, let $D$ be a $(\mu, \epsilon)$-decomposition of $S$. If, each $D_i \in D$ has the trivial partition as its only $(\mu, \epsilon)$-robust partition then the only $(\mu, \epsilon)$-robust partition on $S$ is trivial.
\end{theorem}
% lax end

\begin{proof}   
    We prove the condition (2.) from Theorem~\ref{thm:existence} directly in two steps. Pick a nonempty $S' \subset S$. In the first step, we rule out existence of some $S'$ and in the second, construct the witness $x \in S$ that satisfies Theorem~\ref{thm:existence} (2.).
    
    (\emph{Step 1}): 
    First, we prove that there exists $D_i \in D$ such that both $D_i \cap S'$ and $D_i \cap (S \setminus S')$ are nonempty. For contradiction, assume that $D_i \cap S' = \emptyset$ or $D_i \cap (S \setminus S') = \emptyset$ for all $D_i \in D$. As $S'$ is a proper nonempty subset, there are $x' \in S'$ and $x^* \in S \setminus S'$. Using property (2.) from Definition~\ref{def:decomposition}, construct a sequence $D_{i_1}, \dots, D_{i_p}$ out of elements in $D$ with $x' \in D_{i_1}$ and $x^* \in D_{i_p}$. As $D_{i_1} \subseteq S'$ and $D_{i_p} \cap S' = \emptyset$ by assumption of this step, there is a consecutive pair $D_{i_j}, D_{i_{j + 1}}$  with $j \in [p]$ in the sequence such that $D_{i_j} \cap S' = D_{i_j}$ and $D_{i_{j + 1}} \cap S' = \emptyset$. But, by property (2.) of Definition~\ref{def:decomposition}, $D_{i_j} \cap D_{i_{j + 1}} \neq \emptyset$, leading to contradiction: $\emptyset \neq D_{i_j} \cap D_{i_{j + 1}} = (D_{i_j} \cap S') \cap D_{i_{j + 1}} = D_{i_j} \cap (S' \cap D_{i_{j + 1}}) = \emptyset$.

    (\emph{Step 2}): 
    We construct the witness $x \in S$ in this step. Per Step 1, pick $D_i$ with $D_i \cap S'$ and $D_i \cap (S \setminus S')$ being both nonempty. Then, as $D_i$ does not have nontrivial robust partitions by theorem assumption, by Theorem~\ref{thm:existence}, there exists $x \in D_i$ such that either $x \in D_i \cap S'$ and $(B(x,\mu) \cap D_i) \subseteq \bigcup_{x' \in D_i \setminus S'} (B(x', \epsilon) \cap D_i)$ or $x \in D_i \cap (S \setminus S')$ with $(B(x,\mu) \cap D_i) \subseteq \bigcup_{x' \in D_i \cap S'} (B(x', \epsilon) \cap D_i)$. We pick $x$ as the witness for our proof. If $x \in S'$ then by property (1.) of Definition~\ref{def:decomposition}, $(B(x,\mu) \cap D_i) = B(x, \mu) \subseteq \bigcup_{x' \in D_i \setminus S'} (B(x', \epsilon) \cap D_i) \subseteq \bigcup_{x' \in S \setminus S'} B(x', \epsilon)$. Similarly if $x \in S \setminus S'$, $(B(x,\mu) \cap D_i) = B(x, \mu) \subseteq \bigcup_{x' \in D_i \cap S'} (B(x', \epsilon) \cap D_i) \subseteq \bigcup_{x' \in S'} B(x', \epsilon)$ . As in both cases, condition (2.) from Theorem~\ref{thm:existence} is proven, and $S$ does not have a nontrivial robust partition in this case.
    \qed
\end{proof}

We specialize the result from Theorem~\ref{thm:decomposition} to ordered semimetric spaces.

% lax begin Lax544225.PartialOrderPrecisionBound
\begin{corollary}
    \label{cor:bounds-posets}
    Let $(S, <, d)$ a finite partially ordered semimetric space with $|S| \geq 2$ and $\mu \geq 0, \epsilon > 0$. Let $D := \{D_1, \dots, D_n\}$ be a $(\mu, \epsilon)$-decomposition of $S$ such that all $D_i \in D$ form a totally ordered semimetric space $(D_i, <, d)$ with corresponding $\Delta_i, \delta_i$. If $\epsilon \geq \max_{i \in [n]}(\Delta_i \cdot \delta_i \cdot \mu + \Delta_i)$ then the only $(\mu, \epsilon)$-robust partition of $S$ is trivial.
\end{corollary}
% lax end

\begin{proof}
    First, observe that each $D_i \in D$ does not have a nontrivial $(\mu, \epsilon)$-robust partition by Corollary~\ref{cor:bounds-total}. This implies by Theorem~\ref{thm:decomposition} that every $(\mu, \epsilon)$-robust partition on $S$ is trivial.
\end{proof}

Corollary~\ref{cor:bounds-posets} presents an upper bound for any partially ordered set that has a $(\mu,\epsilon)$-decomposition. Despite our motivation, a general product order does not necessarily admit such decomposition. This means that the Corollary does not easily generalize to vectors of finite precision numbers.

\section{Conclusion}

\paragraph{Discussion.} This work adresses the key limitation of global $\epsilon$-robustness: its stringent requirements on classifiers. As clearly visible in Theorem~\ref{thm:bounds-total}, parameter $\mu$ allows for larger values of $\epsilon$ in totally ordered semimetric spaces such as floating-point numbers (Section~\ref{subsec:existence-order}). $(\mu,\epsilon)$-robustness is proven equivalent in Theorem~\ref{thm:separability} to a coverage condition in parallel to $\epsilon$-robustness being equivalent to separability. This leads to necessary and sufficient conditions for existence of a non-constant $(\mu, \epsilon)$-classifier in general case (Theorem~\ref{thm:existence}) and for totally ordered semimetric spaces (Corollary~\ref{cor:existence-total}). Precision-dependent upper bounds (Corollaries~\ref{cor:bounds-total} and~\ref{cor:bounds-posets}) create connection between underlying arithmetic and existence of non-constant classifiers. By setting $\mu = 0$, one obtains well-known results on $\epsilon$-robustness that were never collected and formally proven in a single place.

\paragraph{Future work.} The next logical step is to design a verification algorithm for $(\mu,\epsilon)$-robustness. This is a significant open research problem, because the $(\mu,\epsilon)$-robustness property is in the class of $\forall\exists$-hyperproperties. While these are studied in general context, neural network verification currently focuses on the domain of hyperproperties with $\forall$ quantifiers only.


%
% ---- Bibliography ----
%
% BibTeX users should specify bibliography style 'splncs04'.
% References will then be sorted and formatted in the correct style.
%
 \bibliographystyle{splncs04}
 \bibliography{sources}
%

\end{document}
